\documentclass[11pt]{article} 
\renewcommand{\thempfootnote}{\arabic{mpfootnote}}

\usepackage[hidelinks]{hyperref}

\usepackage[utf8]{inputenc}

\usepackage[T1]{fontenc}
\usepackage[ttdefault=true]{AnonymousPro}

\usepackage{multicol}
\setlength{\columnsep}{1cm}

\usepackage{abc}
\usepackage{graphicx}
\usepackage{musicography}

\usepackage[margin=0.5in]{geometry} 
\geometry{letterpaper}

\usepackage{fancyhdr} 
\pagestyle{fancy} % options: empty , plain , fancy
\renewcommand{\headrulewidth}{0pt} % customise the layout...
\lhead{}\chead{}\rhead{}
\lfoot{}\cfoot{\thepage}\rfoot{}

\begin{document}

\thispagestyle{empty}

\begin{abc}[name=coverley,options={-O= +c -F single.fmt}]
X:1
T:Sir Roger de Coverley
C:as played by Vivian Williams
C:transcribed by Kitty Steetle
M:9/8
K:G
|: "G"G2G G2e dBG | "D"A2A ABG FED | "G"G2G G2e "D"def | "G"g2G "C"GAG "D"FED :|
|: "G"dBd "C"ece "G"dBG | "D"A2A ABG FED | "G"dBd "C"ece "G"def | "G"g2G "C"GAG "D"FED :| 
|: "G"gag f2e dBG | "Am"A2A ABG "D"FED | "G"gag f2e "D"def | "G"g2G "C"GAG "D"FED :|
\end{abc}
\vskip 0.25in
\noindent\makebox[\textwidth][c]{%
  \begin{minipage}{7.0in}
    \setlength{\parindent}{1.5em}

    The tune, \textit{Sir Roger de Coverley}\footnote{The Haymakers'
      Jig is the Scottish name for the tune. The Haymakers' Jig also
      has a dance in Linscott's \textbf{Folk Songs of Old New England}
      (CT, 1939), p 83, but it is a different dance.}, is believed to
    date back to at least the 17th century, being published in John
    Playford's \textbf{The Dancing Master}, 9th Edition (London,
    1695).\footnote{The entire Dancing Master can be downloaded here:
      \url{https://imslp.org/work/144608}. `Roger of Coverly' can be
      found with dance steps in the 9th edition on page 167.} By 1730,
    this tune appears in America in James Alexander's personal
    collection as `Young Roger'.\footnote{Van Winkle Keller, Kate and
      George A. Fogg. \textbf{Country Dances from Colonial New York;
        James Alexander's Notebook, 1730}. Boston, 2000, p 5, 42 –
      43. James Alexander was born in Scotland in 1691 and moved to
      the Colonies in 1715, eventually becoming Attorney General for
      New York.} The dance associated with the tune survives today in
    America as The Virginia Reel, long claimed to be George
    Washington's favorite
    dance.\footnote{\url{https://www.auchitya.com/virginia-reel-a-romantic-couple-dance/}}
    Perhaps the most famous appearance of the Sir Roger de Coverley
    tune and dance is in Charles Dickens's book, \textbf{A Christmas
      Carol} (London, 1843), during Mr Fezziwig's Christmas party:

    \begin{quotation}
      But the great effect of the evening came after the Roast and
      Boiled, when the fiddler (an artful dog, mind! The sort of man
      who knew his business better than you or I could have told it
      him!)  struck up ``Sir Roger de Coverley.'' Then old Fezziwig
      stood out to dance with Mrs. Fezziwig. Top couple, too; with a
      good stiff piece of work cut out for them; three or four and
      twenty pair of partners; people who were not to be trifled with;
      people who \textit{would} dance, and had no notion of
      walking.\footnote{For Dickens's A Christmas Carol ebook go to:
        \url{https://www.gutenberg.org/files/46/46-h/46-h.htm}}
    \end{quotation}

    Both tune, Sir Roger de Coverley, and dance are performed in the
    1951 film, \textbf{A Christmas Carol}, starting at 26:40, or in
    the YouTube snip below.\footnote{Unfortunately, it does not appear
      in arguably the best \textbf{A Christmas Carol} film, \textbf{A
        Muppet Christmas Carol} (1992).}
    
    Sir Roger de Coverley is notably different from most Old-Time
    tunes in that it is in 9/8, triple meter, instead of the expected
    4/4, duple meter. This has led some to believe it was originally
    Irish, but upon deeper investigation that appears to be
    unsubstantiated.\footnote{\url{https://tunearch.org/wiki/Annotation:Sir\_Roger\_de\_Coverley}
      (I love me some old tunes!)}
    
    Versions of Sir Roger de Coverley are played in Old Time jams, Irish Sessions, English Country Dances,
    and contradances all over North America and Europe.
    
    ---Kitty Steetle

    \bigskip\noindent
    Vivian and Phil Williams: \url{https://youtu.be/y4XB8g4gwLw} \\
    Dance: \url{https://youtu.be/iye3JoD5LUo}  \\
    A Christmas Carol (1951): \url{https://youtu.be/NjwrSy\_YI2Q}
    
    
  \end{minipage}
}


\end{document}
